\section{Resistències en sèrie}

Dues resistències estan \textbf{en sèrie} quan estan una darrere l'altra, en el mateix camí: el corrent no té on escollir i les travessa totes. Pensa en un tub d'aigua amb dues estretors seguides: tota l'aigua passa per les dues.

\begin{center}\cserie{$V_T$}{$R_1$}{$R_2$}\end{center}

\begin{idea}
En un circuit en sèrie:
\begin{enumerate}[label=\textbf{\arabic*.}]
 \item La \textbf{intensitat és la mateixa} a tot arreu: \quad $I_T = I_1 = I_2 = \cdots$
 \item Les \textbf{tensions se sumen}: \quad $V_T = V_1 + V_2 + \cdots$
 \item Les \textbf{resistències se sumen}: \quad $R_T = R_1 + R_2 + \cdots$
\end{enumerate}
\end{idea}

\begin{metode}{calcular un circuit en sèrie}
\begin{itemize}
 \item[\textbf{\ding{172}}] Suma les resistències: $R_T = R_1 + R_2 + \cdots$
 \item[\textbf{\ding{173}}] Calcula la intensitat total: $I_T = V_T / R_T$. Aquesta és la intensitat de \emph{cada} resistència.
 \item[\textbf{\ding{174}}] Calcula la tensió de cada resistència: $V_i = I \cdot R_i$.
 \item[\textbf{\ding{175}}] Calcula la potència de cada resistència: $P_i = V_i \cdot I$.
 \item[\textbf{\ding{176}}] \textbf{Comprova}: les tensions sumen $V_T$? les potències sumen $V_T \cdot I_T$?
\end{itemize}
\end{metode}

\begin{exemple}{dues resistències en sèrie}
Una pila de \SI{12}{V} alimenta dues resistències en sèrie: $R_1 = \SI{100}{\ohm}$ i $R_2 = \SI{200}{\ohm}$.
\begin{center}\cserie{\SI{12}{V}}{$R_1=\SI{100}{\ohm}$}{$R_2=\SI{200}{\ohm}$}\end{center}
\begin{itemize}
 \item $R_T = 100 + 200 = \SI{300}{\ohm}$
 \item $I_T = V_T / R_T = 12 / 300 = \SI{0,04}{A}$ \ (la mateixa per a $R_1$ i $R_2$)
 \item $V_1 = I \cdot R_1 = 0{,}04 \cdot 100 = \SI{4}{V}$ \qquad $V_2 = 0{,}04 \cdot 200 = \SI{8}{V}$
 \item $P_1 = 4 \cdot 0{,}04 = \SI{0,16}{W}$ \qquad $P_2 = 8 \cdot 0{,}04 = \SI{0,32}{W}$ \qquad $P_T = 12 \cdot 0{,}04 = \SI{0,48}{W}$
 \item \textbf{Comprovació:} $4 + 8 = 12$ \checkmark \qquad $0{,}16 + 0{,}32 = 0{,}48$ \checkmark
\end{itemize}
\begin{taulamag}
$R_1$ & 4 & 0{,}04 & 100 & 0{,}16\\
$R_2$ & 8 & 0{,}04 & 200 & 0{,}32\\
\midrule
Total & 12 & 0{,}04 & 300 & 0{,}48\\
\end{taulamag}
\end{exemple}

\begin{idea}
\textbf{Un truc útil: el divisor de tensió.} En sèrie, la tensió es reparteix \emph{en proporció a la resistència}. La resistència més gran se'n queda la part més gran:
\[ V_i = V_T \cdot \frac{R_i}{R_T} \]
A l'exemple: $V_1 = 12 \cdot \dfrac{100}{300} = \SI{4}{V}$.
\end{idea}

\begin{atencio}
Sumar les intensitats en un circuit en sèrie. \textbf{No!} La intensitat és \emph{la mateixa} a totes les resistències; el que se suma són les \textbf{tensions} i les \textbf{resistències}.
\end{atencio}

\begin{guiada}{pila de \SI{9}{V} amb $R_1 = \SI{300}{\ohm}$ i $R_2 = \SI{600}{\ohm}$ en sèrie}
\pas{1}{Resistència total:}{$R_T = $ \vspace{0.7cm}}
\pas{2}{Intensitat total:}{\vspace{0.7cm}}
\pas{3}{Tensions de cada resistència:}{\vspace{0.7cm}}
\begin{taulamag}
$R_1$ & & & 300 &\\
$R_2$ & & & 600 &\\
\midrule
Total & 9 & & &\\
\end{taulamag}
\resposta{0pt}{$R_T = 900\,\Omega$; $I = 9/900 = \SI{0,01}{A}$; $V_1 = 0{,}01 \cdot 300 = \SI{3}{V}$; $V_2 = 0{,}01 \cdot 600 = \SI{6}{V}$; $P_1 = \SI{0,03}{W}$, $P_2 = \SI{0,06}{W}$, $P_T = \SI{0,09}{W}$.
\begin{taulamag}
$R_1$ & 3 & 0{,}01 & 300 & 0{,}03\\
$R_2$ & 6 & 0{,}01 & 600 & 0{,}06\\
\midrule
Total & 9 & 0{,}01 & 900 & 0{,}09\\
\end{taulamag}}
\end{guiada}

\subsection*{Exercicis}

\begin{exercici}{1}
Calcula la resistència total de:\quad
\textbf{a)} \SI{100}{\ohm} i \SI{220}{\ohm} en sèrie.\qquad
\textbf{b)} \SI{1}{\kilo\ohm}, \SI{2,2}{\kilo\ohm} i \SI{4,7}{\kilo\ohm} en sèrie.
\resposta{2cm}{a) $100 + 220 = \SI{320}{\ohm}$. \quad b) $1 + 2{,}2 + 4{,}7 = \SI{7,9}{\kilo\ohm}$.}
\end{exercici}

\begin{exercici}[20]{1}
Completa la taula del circuit següent.
\begin{center}\cserie{\SI{9}{V}}{$R_1=\SI{100}{\ohm}$}{$R_2=\SI{200}{\ohm}$}\end{center}
\begin{taulamag}
$R_1$ & & & 100 &\\
$R_2$ & & & 200 &\\
\midrule
Total & 9 & & &\\
\end{taulamag}
\resposta{0pt}{$R_T = \SI{300}{\ohm}$, $I = \SI{0,03}{A}$.
\begin{taulamag}
$R_1$ & 3 & 0{,}03 & 100 & 0{,}09\\
$R_2$ & 6 & 0{,}03 & 200 & 0{,}18\\
\midrule
Total & 9 & 0{,}03 & 300 & 0{,}27\\
\end{taulamag}}
\end{exercici}

\begin{exercici}[20]{2}
Completa la taula del circuit següent.
\begin{center}\cserie{\SI{24}{V}}{$R_1=\SI{300}{\ohm}$}{$R_2=\SI{500}{\ohm}$}\end{center}
\begin{taulamag}
$R_1$ & & & 300 &\\
$R_2$ & & & 500 &\\
\midrule
Total & 24 & & &\\
\end{taulamag}
\resposta{0pt}{$R_T = \SI{800}{\ohm}$, $I = \SI{0,03}{A}$.
\begin{taulamag}
$R_1$ & 9 & 0{,}03 & 300 & 0{,}27\\
$R_2$ & 15 & 0{,}03 & 500 & 0{,}45\\
\midrule
Total & 24 & 0{,}03 & 800 & 0{,}72\\
\end{taulamag}}
\end{exercici}

\begin{exercici}[20]{2}
Completa la taula del circuit següent, amb tres resistències en sèrie.
\begin{center}\cserietres{\SI{12}{V}}{$R_1=\SI{100}{\ohm}$}{$R_2=\SI{200}{\ohm}$}{$R_3=\SI{300}{\ohm}$}\end{center}
\begin{taulamag}
$R_1$ & & & 100 &\\
$R_2$ & & & 200 &\\
$R_3$ & & & 300 &\\
\midrule
Total & 12 & & &\\
\end{taulamag}
\resposta{0pt}{$R_T = \SI{600}{\ohm}$, $I = \SI{0,02}{A}$.
\begin{taulamag}
$R_1$ & 2 & 0{,}02 & 100 & 0{,}04\\
$R_2$ & 4 & 0{,}02 & 200 & 0{,}08\\
$R_3$ & 6 & 0{,}02 & 300 & 0{,}12\\
\midrule
Total & 12 & 0{,}02 & 600 & 0{,}24\\
\end{taulamag}}
\end{exercici}

\begin{exercici}{2}
Dues resistències en sèrie, $R_1 = \SI{150}{\ohm}$ i una $R_2$ desconeguda, estan connectades a \SI{10}{V}. Per elles hi passa un corrent de \SI{0,05}{A}. Calcula $R_2$ i la tensió de cada resistència.
\resposta{3cm}{$R_T = V/I = 10/0{,}05 = \SI{200}{\ohm}$, per tant $R_2 = 200 - 150 = \SI{50}{\ohm}$. $V_1 = 0{,}05 \cdot 150 = \SI{7,5}{V}$, $V_2 = 0{,}05 \cdot 50 = \SI{2,5}{V}$. Comprovació: $7{,}5 + 2{,}5 = 10$ \checkmark}
\end{exercici}

\begin{exercici}{3}
Volem fer un \textbf{divisor de tensió}: una font de \SI{5}{V} alimenta $R_1 = \SI{300}{\ohm}$ i una $R_2$ en sèrie, i volem que entre els extrems de $R_2$ hi hagi exactament \SI{2}{V}. Quant ha de valer $R_2$?
\resposta{3.5cm}{$V_1 = 5 - 2 = \SI{3}{V}$. $I = V_1 / R_1 = 3/300 = \SI{0,01}{A}$. $R_2 = V_2 / I = 2 / 0{,}01 = \SI{200}{\ohm}$. Comprovació amb el divisor: $V_2 = 5 \cdot 200/500 = \SI{2}{V}$ \checkmark}
\end{exercici}

\begin{exercici}{3}
Una cadena de llums de Nadal té \num{20} bombetes iguals en sèrie i es connecta a \SI{230}{V}.
\textbf{a)} Quina tensió suporta cada bombeta?\quad \textbf{b)} Què passa amb la resta si una bombeta es fon? Per què?
\resposta{3cm}{a) $230 / 20 = \SI{11,5}{V}$ cada bombeta. b) S'apaguen totes: en fondre's una, el circuit s'obre i el corrent no pot circular enlloc.}
\end{exercici}